\documentclass[a4paper,12pt]{article}

% 基础宏包配置
\usepackage{ctex}      % 中文支持
\usepackage{geometry}  % 页面布局
\usepackage{titlesec}  % 标题格式
\usepackage{fancyhdr}  % 页眉页脚
\usepackage{xcolor}    % 颜色支持
\usepackage{setspace}  % 行间距设置

% 页面边距设置
\geometry{left=2.5cm, right=2.5cm, top=3cm, bottom=3cm}
\onehalfspacing % 1.5倍行距，符合正式公文规范

% 颜色定义
\definecolor{headerBlue}{RGB}{0, 51, 102}
\definecolor{subHeaderGray}{RGB}{80, 80, 80}

% 标题格式定制
\titleformat{\section}
{\normalfont\Large\bfseries\color{headerBlue}}{\thesection}{1em}{}
\titleformat{\subsection}
{\normalfont\large\bfseries\color{subHeaderGray}}{\thesubsection}{1em}{}

% 页眉页脚
\pagestyle{fancy}
\fancyhf{}
\lhead{\small \textbf{2024年度ESG报告}}
\rhead{\small 责任管理与绩效披露}
\cfoot{\thepage}

\title{\textbf{\huge 2024年度环境、社会及治理（ESG）报告}\\ \vspace{0.5em} \large 管理层讨论与分析及可持续发展绩效}
\author{本公司董事会}
\date{2025年3月}

\begin{document}

\maketitle

\section{治理篇：稳健经营与合规管理}

2024年，本公司坚持“三直四化”（电动化、智能网联化、高端化、国际化）战略导向，在复杂多变的全球经济环境中保持了稳健的发展态势。公司深刻理解ESG治理对企业长期价值的影响，致力于将可持续发展理念融入公司治理、战略决策及日常运营之中。

\subsection{经济绩效与股东权益}

报告期内，受益于海外高端市场的拓展及国内新能源客车市场的复苏，公司实现了经营业绩的显著增长。2024年度，公司实现营业收入372.18亿元，同比增长37.63\%；实现归属于母公司股东的净利润41.16亿元。

公司高度重视对投资者的合理回报，坚持实施积极、持续、稳定的利润分配政策。为切实回馈股东，公司在报告期内实施了两次现金分红，全年累计派发现金股利44.28亿元，分红总额超过当期净利润，充分体现了公司现金流的稳健性及对股东权益的保障。

\subsection{合规运营与风险控制}

合规经营是公司可持续发展的基石。公司持续完善内部控制体系，强化合规管理效能。
\begin{itemize}
    \item \textbf{信息披露管理}：公司严格遵循监管要求，提升信息披露的透明度与规范性。报告期内，公司连续第13年获得证券交易所信息披露工作“A级”评价。
    \item \textbf{合规体系认证}：2024年，公司顺利通过ISO 37301合规管理体系认证，标志着公司的合规管理架构、运行机制及保障措施已达到国际标准要求。
\end{itemize}

\subsection{供应链可持续管理}

公司致力于构建绿色、廉洁、高效的供应链生态。截至2024年末，公司拥有合作供应商530家。公司积极发挥链主企业的辐射带动作用，支持属地化采购，其中河南本地供应商占比达到20\%，有效促进了区域装备制造业的协同发展。

在商业道德方面，公司严格执行反腐败与反商业贿赂政策。报告期内，公司已与超过900家供应商（含潜在供应商）签署廉洁协议，明确双方在商业交往中的廉洁责任，确保供应链的透明与公正。

\section{社会篇：创新驱动与共享价值}

公司坚持技术创新作为核心驱动力，推动商用车行业的绿色化与智能化转型。同时，公司坚持“以人为本”，切实保障员工权益，并积极投身社会公益事业，与利益相关方共享发展成果。

\subsection{研发创新与技术成果}

公司持续保持高强度的研发投入，以巩固在新能源与智能网联领域的技术领先优势。2024年，公司研发投入总额达17.88亿元，占营业收入比例为4.80\%。公司拥有一支高素质的研发团队，研发人员总数达3,505人，其中包含博士16人。

在知识产权方面，报告期内公司新增授权专利176项，其中发明专利114项；累计参与制定国家及行业标准317项，持续引领行业技术规范的发展。

\subsubsection{关键核心技术进展}

2024年，公司在车辆软硬件架构及核心零部件领域取得多项突破：
\begin{itemize}
    \item \textbf{软件定义汽车}：自主研发“YOS车机操作系统”，实现了软硬件解耦。该系统支持全栈OTA升级，大幅提升了车辆全生命周期的功能迭代能力。
    \item \textbf{电子电气架构}：全面应用“C架构”（中央计算+区域控制），通过高集成度设计，有效减少线束冗余，提升了整车的通信效率与轻量化水平。
    \item \textbf{高效电驱动系统}：开发并应用了多合一动力域控制器，采用碳化硅（SiC）功率模块。相较于传统硅基模块，其开关损耗降低50\%-70\%，显著提升了整车能源利用效率。
    \item \textbf{热管理技术}：推出了多源超低温热泵系统，通过高效回收环境热能及电机余热，在低温工况下采暖节能效果达到30\%，有效改善了新能源车辆冬季续航性能。
\end{itemize}

此外，在智能网联运营方面，公司L4级自动驾驶巴士已累计安全运营超过1,700万公里，并成功获批国家首批智能网联汽车准入试点，验证了技术的成熟度与商业化潜力。

\subsection{人力资源发展与安全生产}

截至报告期末，集团员工总数为16,707人，其中女性员工占比10.02\%。公司建立了完善的职业健康安全管理体系（ISO 45001），坚持“安全第一，预防为主”的方针。

报告期内，公司在安全生产领域投入力度持续加大。其中，劳动防护用品投入5,500余万元，用于提升一线员工的个人防护水平；生产现场环境改善投入2,200余万元，用于除尘、降噪及通风设施的升级改造；安全教育培训投入180余万元。2024年，公司累计开展安全培训7.9万人次，全年未发生新增职业病病例。

\subsection{企业公民与社区贡献}

公司积极履行社会责任，2024年对外捐赠及公益项目总投入超过4,000万元。

在品牌公益方面，“儿童交通安全公益行”项目持续开展，足迹遍布全国13个省市。该项目通过模拟车辆盲区等互动体验形式，向近万名儿童普及了交通安全知识。在海外市场，公司在南美某国实施了“零碳森林”项目，针对当地生态需求种植耐旱树种，累计植树36,000余棵，在助力全球碳减排的同时，改善了当地社区的生态环境。

\section{环境篇：低碳转型与绿色制造}

公司严格遵守相关环境法律法规，将绿色低碳理念贯穿于产品设计、原材料采购、生产制造及物流运输的全过程。

\subsection{能源管理与温室气体排放}

公司建立了完善的碳排放数据监测与核算体系。2024年，公司综合能源消费量为34,560.94吨标准煤。

在温室气体排放方面，报告期内范围1（直接排放）排放量为57,330.53吨二氧化碳当量；范围2（能源间接排放）排放量为122,786.63吨二氧化碳当量。经测算，公司温室气体排放强度为0.07吨二氧化碳/万元营收。

为降低碳排放强度，公司实施了多项节能降碳措施：
\begin{itemize}
    \item \textbf{节能技改}：全年实施重点节能改造项目7项，通过设备能效提升与工艺优化，实现年节电1,258兆瓦时，年节气11.34万立方米。
    \item \textbf{余热回收}：大力推广余热回收技术，全年回收利用工业余热约29,306吉焦，替代了部分化石能源消耗。
    \item \textbf{清洁能源利用}：公司积极优化能源结构，新能源智造基地光伏发电系统全年发电量达155.68万千瓦时，折合减碳约835吨；此外，公司全年采购绿色电力13,580兆瓦时。
\end{itemize}

\subsection{资源循环与废弃物合规处置}

公司坚持循环经济原则，提升资源利用效率。水资源管理方面，2024年新取水量为141.25万吨，通过中水回用系统，重复用水量达到6,465.36万吨，重复用水率高达97.8\%。

废弃物管理方面，公司严格执行分类收集与合规处置。报告期内，一般工业固废中回收利用量为61,485.76吨；危险废物产生量为4,110.84吨，全部委托具备相应资质的单位进行无害化处置，处置合规率保持100\%。

此外，公司在物流环节积极推广绿色包装应用。目前，国产零部件采购环节的周转包装（可循环利用包装）占比已达到89\%，显著减少了物流过程中的一次性包装废弃物。

\section{未来展望}

展望未来，本公司将继续深化ESG管理实践，依托“三直四化”战略，持续探索氢能技术应用、供应链碳足迹管理等前沿议题。我们将坚持技术创新与责任担当并重，携手各方利益相关者，共同推动交通运输行业的绿色、智能与可持续发展。

\vspace{2cm}

\end{document}